\documentclass[a4paper, 12pt]{article}
\usepackage{xeCJK}
\usepackage{amsmath}
\usepackage{fancyhdr}
\usepackage{geometry}
\usepackage{tikz}
\begin{document}
\section{传送带问题}
\subsection*{水平传送带}
\subsubsection*{同向运动}
传送带速度为$v_0$，物体速度为$v$。\\
\textbf{1. $v>v_0$}\\
\begin{tikzpicture}
\draw[->] (0,0) -- (6,0) node [below] {$t$};
\draw[->] (0,0) node (v1) {} -- (0,5) node [right] {$v$};
\draw [dashed, red] (0,2.5) -- (6,2.5) node[pos=0, left] {$v_0$};
\draw (4,2.5) -- (1.5,2.5) -- (0,4) node [pos=1, left] {$v_1$};
\end{tikzpicture}\\
\textbf{2. $v<v_0$}\\
\begin{tikzpicture}
\draw[->] (0,0) -- (6,0) node [below] {$t$};
\draw[->] (0,0) node (v1) {} -- (0,5) node [right] {$v$};
\draw [dashed, red] (0,2.5) -- (6,2.5) node[pos=0, left] {$v_0$};
\draw[blue] (0,0) -- (3,2.5) -- (6,2.5);
\end{tikzpicture}
\subsubsection*{反向运动}
\begin{tikzpicture}
\draw[->] (0,0) -- (6,0) node[below] {$t$};
\draw[->] (0,-2) -- (0,2) node[left] {$v$};
\draw[dashed, red] (0,1.5) -- (6,1.5) node[pos=0, left] {$v_0$};
\draw (4.5,1.5) -- (3,1.5) -- (0,-1) node [left] {$v_1$};
\end{tikzpicture}
\subsection*{倾斜传送带}
传送带速度为$v_0$，物体速度为$v$（如果有的话）。定义$a_1=g(\sin\theta+\mu\cos\theta)$，
$a_2=g|\sin\theta-\mu\cos\theta|$。\\
\subsubsection*{底部，传送带向上}
\textbf{1. $\mu > \tan\theta$}\\
物块有初速度，大于传送带速度，\\
\begin{tikzpicture}
\draw[->] (0,0) -- (6,0) node[below] {$t$};
\draw[->] (0,0) -- (0,3) node[left] {$v$};
\draw[dashed, red] (0,1.5) -- (6,1.5) node[pos=0, left] {$v_0$};
\draw (4.5,1.5) -- (2.5,1.5) -- (0,2.5) node[pos=1, left] {$v_1$} node[pos=0, above, midway] {$a_1$};
\end{tikzpicture}
物块初速度小于传送带速度，\\
\begin{tikzpicture}
\draw[->] (0,0) -- (6,0) node[below] {$t$};
\draw[->] (0,0) -- (0,3) node[left] {$v$};
\draw[dashed, red] (0,1.5) -- (6,1.5) node[pos=0, left] {$v_0$};
\draw (4.5,1.5) -- (2.5,1.5) -- (0,0.5) node[pos=1, left] {$v_1$} node[pos=0, above, midway] {$a_2$};
\end{tikzpicture}
\textbf{2. $\mu < \tan\theta$}\\
物块必须要有初速度，\\

\subsubsection*{顶部，传送带向下}
\textbf{1. $\mu > \tan\theta$}\\
物块有初速度，并且初速度大于传送带速度，\\
物块初速度小于传送带速度，\\
\textbf{2. $\mu < \tan\theta$}\\
\section{板块问题}
物体初速度定义为$v_1$，木板初速度定义为$v_2$，物体与木板之间摩擦因数定义为$\mu_1$，木板
与地面之间摩擦因数定义为$\mu_2$，系统受到的外力定义为$F$。\\
\subsection*{地面光滑}
\textbf{1. 物体有初速度}\\
\textbf{2. 木板有初速度}\\
\textbf{3. 物体有外力}\\
\textbf{4. 木板有外力}\\
\subsection*{地面粗糙}
\textbf{1. 物体有初速度}\\
\textbf{2. 木板有初速度}\\
\textbf{3. 物体有外力}\\
\textbf{4. 木板有外力}\\
\end{document}